\documentclass[11pt]{article}

\usepackage[margin=1.1in]{geometry}
\usepackage{amsmath,amssymb,amsthm,mathtools}
\usepackage{enumitem}
\usepackage[hidelinks]{hyperref}

\newtheorem{theorem}{Theorem}[section]
\newtheorem{lemma}[theorem]{Lemma}
\newtheorem{corollary}[theorem]{Corollary}
\newtheorem{proposition}[theorem]{Proposition}
\theoremstyle{definition}
\newtheorem{definition}[theorem]{Definition}
\newtheorem{remark}[theorem]{Remark}
\newcommand{\docversion}{20}
\newcommand{\rad}{\operatorname{rad}}
\newcommand{\lcm}{\operatorname{lcm}}
\newcommand{\ord}{\operatorname{ord}}
\newcommand{\QR}{\mathbb{Q}}
\newcommand{\Z}{\mathbb{Z}}
\newcommand{\N}{\mathbb{N}}
\newcommand{\Pmax}{P^+}

\title{A Small-Index Reduction for Prime-Complete Products of Two Consecutive Integers (Version \docversion)}
\author{Ken Clements}
\date{\today}

\begin{document}
\maketitle

\begin{abstract}
Let
\[
        P_r=p_r^\#=\prod_{i=1}^r p_i
\]
be the $r$-th primorial.  A product $m(m+1)$ is called prime-complete of order $r$ if
\[
        \rad(m(m+1))=P_r.
\]
This note proves an unconditional reduction theorem for any such product in the tail range.  
If $D$ is the squarefree part of $m(m+1)$, then $D\mid P_r$ and $x=2m+1$ lies on the Pell branch
\[
        X^2-DY^2=1.
\]
Writing the positive solutions of this Pell equation as
\[
        X_n+Y_n\sqrt D=(X_1+Y_1\sqrt D)^n,
\]
the corresponding index $n$ must be small:
\[
        n \le p_r+1.
\]
Moreover every prime $p\le p_r$ not dividing $D$ must divide $Y_n$, and hence the least 
index $\rho_p(D)$ for which $p\mid Y_{\rho_p(D)}$ must divide $n$.  
Therefore the least-common-multiple obstruction
\[
        R_r(D)=\lcm_{\substack{3\le p\le p_r\\ p\nmid D}} \rho_p(D)
\]
must satisfy
\[
        R_r(D)\mid n\le p_r+1.
\]
Consequently a branch $D$ with $R_r(D)>p_r+1$ cannot support a prime-complete product of order $r$.

The proof uses only the Pell reduction, the divisibility properties of Lucas sequences, 
and the Bilu--Hanrot--Voutier primitive divisor theorem.  It does not use the $abc$ 
conjecture, any bound for the largest pair of consecutive $p_r$-smooth integers, 
or any lower bound for the first CRT root.
\end{abstract}

\section{Introduction}

The problem considered here is the termination of prime-complete products of two consecutive integers.  
We say that $m(m+1)$ is prime-complete of order $r$ when its radical is exactly the $r$-th primorial:
\[
        \rad(m(m+1))=P_r=p_r^\#.
\]
Equivalently, the set of distinct prime factors of $m(m+1)$ is precisely
\[
        \{2,3,5,\ldots,p_r\}.
\]

One possible route to termination is to compare two global objects: a smooth ceiling for 
consecutive $p_r$-smooth pairs and a discrete radical--CRT floor.  The reduction proved 
here takes a different route.  It does not attempt to bound all consecutive smooth pairs.  
Instead, it keeps the prime-complete condition and asks what must happen on the Pell 
branch attached to a hypothetical prime-complete product.

The result is a small-index reduction.  Any prime-complete product of order $r$ gives 
a Pell solution indexed by some $n$.  The Bilu--Hanrot--Voutier primitive divisor 
theorem bounds that index by $p_r+1$ in the tail range, while rank-of-apparition 
divisibility forces a large least common multiple to divide the same index.  
Thus, for each squarefree branch $D\mid P_r$, a simple necessary condition emerges:
\[
        R_r(D) \le p_r+1.
\]
Branches violating this inequality are eliminated unconditionally.

This note is deliberately limited in scope.  It does not prove that no prime-complete 
product exists in the tail.  Rather, it proves that any tail counterexample must 
live in a sharply constrained small-index core.

\section{Definitions and notation}

Let $p_r$ denote the $r$-th prime and set
\[
        P_r=p_r^\#=\prod_{i=1}^r p_i.
\]
For a positive integer $N$, let $\rad(N)$ be the product of the distinct prime 
factors of $N$, and let $P^+(N)$ denote its largest prime factor.

\begin{definition}[Prime-complete consecutive product]
A product $m(m+1)$ is \emph{prime-complete of order $r$} if
\[
        \rad(m(m+1))=P_r.
\]
\end{definition}

If $N=m(m+1)$ is prime-complete of order $r$, then all prime factors of $N$ 
are at most $p_r$, and every prime at most $p_r$ occurs at least once.

\begin{definition}[Squarefree branch]
Let $N=m(m+1)$, and write
\[
        N=Dz^2
\]
with $D$ squarefree and $z\in\Z_{>0}$.  We call $D$ the \emph{Pell branch} 
or \emph{squarefree branch} attached to $m(m+1)$.
\end{definition}

If $m(m+1)$ is prime-complete of order $r$, then
\[
        D \mid P_r.
\]
Moreover $D>1$, since $m(m+1)$ cannot be a positive square for $m \ge 1$.

For a squarefree integer $D>1$, let
\[
        \eta_D=X_1+Y_1\sqrt D
\]
be the fundamental positive solution of
\[
        X^2-DY^2=1.
\]
The positive solutions are indexed by $n\ge1$ via
\[
        X_n+Y_n\sqrt D=\eta_D^n.
\]
The sequence $(Y_n)$ is a Lucas sequence up to the fixed factor $Y_1$:
\[
        U_n(D):=\frac{Y_n}{Y_1}
\]
is an integer Lucas sequence with real roots.

\section{The Pell reduction}

\begin{lemma}[Pell attachment]\label{lem:pell-attachment}
Let $m(m+1)$ be prime-complete of order $r$.  Write
\[
        m(m+1)=Dz^2
\]
with $D$ squarefree.  Then $D\mid P_r$, and
\[
        X=2m+1,\qquad Y=2z
\]
satisfy
\[
        X^2-DY^2=1.
\]
Consequently, for some $n\ge1$,
\[
        X+Y\sqrt D=(X_1+Y_1\sqrt D)^n.
\]
Furthermore $Y_n=Y$ is $p_r$-smooth.
\end{lemma}

\begin{proof}
Since $m(m+1)=Dz^2$, we have
\[
        (2m+1)^2-1=4m(m+1)=4Dz^2.
\]
Thus, with $X=2m+1$ and $Y=2z$,
\[
        X^2-DY^2=1.
\]
Because $\rad(m(m+1))=P_r$, the squarefree part $D$ divides $P_r$.  Also $z$ 
has no prime factor exceeding $p_r$, so $Y=2z$ has no prime factor exceeding 
$p_r$.  Hence $Y$ is $p_r$-smooth.

Every positive solution of $X^2-DY^2=1$ is a positive power of the fundamental 
solution, so the asserted index $n$ exists.
\end{proof}

\begin{definition}[Complementary product]
For $D\mid P_r$, define the complementary product
\[
        B_r(D)=\frac{P_r}{D}.
\]
\end{definition}

\begin{lemma}[Complementary support]\label{lem:complement-support}
Let $m(m+1)$ be prime-complete of order $r$, and let $D$ be its squarefree branch.  
If the attached Pell solution has $Y=Y_n$, then
\[
        B_r(D)\mid Y_n.
\]
In particular, for every prime $p\le p_r$ with $p\nmid D$, one has
\[
        p\mid Y_n.
\]
\end{lemma}

\begin{proof}
Write $m(m+1)=Dz^2$.  Let $p\le p_r$ and suppose $p\nmid D$.  Since $m(m+1)$ is 
prime-complete, $p$ divides $m(m+1)$.  Because $p\nmid D$ and $D$ is the squarefree 
part, the exponent of $p$ in $m(m+1)$ is positive and even.  Hence $p\mid z$, and 
therefore $p\mid Y=2z$.  This proves that every prime dividing $B_r(D)$ divides 
$Y_n$, so $B_r(D)\mid Y_n$.
\end{proof}

\begin{remark}[Parity]
The Pell equation contains solutions that do not necessarily correspond to an 
integer $m$, since an actual consecutive product requires $X$ odd and $Y$ even.  
The reduction in this note is only a necessary condition.  It is harmless, and 
slightly stronger, to allow all Pell terms in the obstruction.  Parity can be 
imposed as an additional filter in computations, but it is not needed for 
the theorem below.
\end{remark}

\section{Primitive divisors and the index cap}

The first jaw of the reduction is a primitive-divisor bound for the Lucas sequence
\[
        U_n(D)=\frac{Y_n}{Y_1}.
\]
We use the following standard consequence of the Bilu--Hanrot--Voutier primitive 
divisor theorem for Lucas and Lehmer sequences.

\begin{lemma}[Primitive divisor bound]\label{lem:primitive}
Let $D>1$ be squarefree, and let $U_n(D)=Y_n/Y_1$ be the Lucas sequence arising 
from the Pell equation $X^2-DY^2=1$.  If $n>30$, then $U_n(D)$ has a primitive 
prime divisor $q$.  Moreover $q$ is coprime to $2D$ and satisfies
\[
        n\mid q-\left(\frac{D}{q}\right),
\]
where $\left(\frac{D}{q}\right)$ is the Legendre symbol.  Consequently
\[
        q\equiv \pm 1 \pmod n
\]
and hence
\[
        q\ge n-1.
\]
\end{lemma}

\begin{proof}
The existence of a primitive prime divisor for $n>30$ is the Bilu--Hanrot--Voutier 
theorem applied to the nondegenerate Lucas sequence associated to the conjugate pair
\[
        \eta_D=X_1+Y_1\sqrt D,
        \qquad
        \eta_D^{-1}=X_1-Y_1\sqrt D.
\]
For a primitive divisor $q$ of $U_n(D)$, the rank of apparition of $q$ in the sequence 
is exactly $n$.  The usual rank divisibility for Pell--Lucas sequences gives
\[
        n\mid q-\left(\frac{D}{q}\right).
\]
Since the Legendre symbol is $\pm 1$, this gives $q\equiv\pm 1 \pmod n$, and therefore $q \ge n-1$.
\end{proof}

\begin{proposition}[Small-index cap]\label{prop:index-cap}
Let $r \ge 10$.  Suppose a Pell term $Y_n$ arising from a branch $D \mid P_r$ is $p_r$-smooth.  Then
\[
        n \le p_r+1.
\]
\end{proposition}

\begin{proof}
If $n \le 30$, then the conclusion holds for $r\ge10$, since $p_{10}=29$ and therefore $p_r+1\ge30$.

Suppose $n>30$.  By Lemma~\ref{lem:primitive}, $U_n(D)$ has a primitive prime divisor 
$q$ satisfying $q\ge n-1$.  Since $U_n(D)$ divides $Y_n$, and $Y_n$ is $p_r$-smooth, 
we have $q \le p_r$.  Therefore
\[
        n-1 \le p_r,
\]
which is equivalent to
\[
        n \le p_r+1.
\]
\end{proof}

\section{Ranks of apparition and support-completeness}

The second jaw of the reduction is the elementary rank-of-apparition condition.

\begin{definition}[Rank of apparition]\label{def:rank}
Let $D>1$ be squarefree and let $(Y_n)$ be the $Y$-sequence of $X^2-DY^2=1$.  
For an odd prime $p\nmid D$, define
\[
        \rho_p(D)=\min\{n\ge1:p\mid Y_n\}.
\]
This minimum exists because the Pell unit has finite order in the norm-one 
subgroup modulo $p$.
\end{definition}

\begin{lemma}[Rank divisibility]\label{lem:rank-divisibility}
Let $D>1$ be squarefree and let $p$ be an odd prime with $p \nmid D$.  Then
\[
        p\mid Y_n
\]
if and only if
\[
        \rho_p(D)\mid n.
\]
\end{lemma}

\begin{proof}
Modulo $p$, work in the quadratic algebra $\mathbb F_p(\sqrt D)$ if $D$ is a 
quadratic nonresidue and in $\mathbb F_p\times\mathbb F_p$ if $D$ is a residue.  The Pell unit
\[
        \eta_D=X_1+Y_1\sqrt D
\]
has norm $1$.  The condition $p\mid Y_n$ is exactly the condition that
\[
        \eta_D^n=X_n+Y_n\sqrt D
\]
has zero $\sqrt D$-coefficient modulo $p$, equivalently that $\eta_D^n$ lies in the 
rational subfield modulo $p$.  Since the norm is 1, this rational element is $\pm 1$.  
Thus the set of $n$ for which $p\mid Y_n$ is a nonzero subgroup of $\Z$ under addition, 
hence is precisely the set of multiples of its least positive element, $\rho_p(D)$.
\end{proof}

\begin{definition}[Rank lcm obstruction]
Let $D\mid P_r$ be squarefree.  Define
\[
        R_r(D)
        =
        \lcm_{\substack{3\le p\le p_r\\ p\nmid D}} \rho_p(D).
\]
The prime $2$ is omitted only to avoid parity exceptions in the formal statement; 
including a valid parity rank for 2 can only strengthen the computational obstruction.
\end{definition}

\begin{lemma}[Complementary rank condition]\label{lem:rank-condition}
Let $m(m+1)$ be prime-complete of order $r$, and let $D$ be its squarefree branch.  
If the attached Pell term is $Y_n$, then
\[
        R_r(D)\mid n.
\]
\end{lemma}

\begin{proof}
By Lemma~\ref{lem:complement-support}, every odd prime $p\le p_r$ with $p\nmid D$ 
divides $Y_n$.  By Lemma~\ref{lem:rank-divisibility}, this implies
\[
        \rho_p(D)\mid n
\]
for every such $p$.  Taking the least common multiple over all odd complementary primes gives
\[
        R_r(D)\mid n.
\]
\end{proof}

\section{The reduction theorem}

We can now state and prove the main reduction.

\begin{theorem}[Small-index reduction theorem]\label{thm:reduction}
Let $r\ge 38$.  Suppose that $m(m+1)$ is prime-complete of order $r$.  Let $D$ be the 
squarefree part of $m(m+1)$, so that $D\mid P_r$, and let
\[
        X_n+Y_n\sqrt D=(X_1+Y_1\sqrt D)^n
\]
be the Pell solution attached to $m$ by
\[
        X_n=2m+1.
\]
Then all of the following hold:
\begin{enumerate}[label=\textup{(\roman*)}]
    \item $Y_n$ is $p_r$-smooth;
    \item $B_r(D)=P_r/D$ divides $Y_n$;
    \item $n\le p_r+1$;
    \item $R_r(D)\mid n$.
\end{enumerate}
In particular, a necessary condition for the branch $D$ to support a prime-complete 
product of order $r$ is
\[
        R_r(D)\le p_r+1.
\]
\end{theorem}

\begin{proof}
Assertions (i) and (ii) are Lemmas~\ref{lem:pell-attachment} and~\ref{lem:complement-support}.  
Since $r \ge 38 \ge 10$, Proposition~\ref{prop:index-cap} applies and gives
\[
        n \le p_r+1,
\]
which proves (iii).  Finally, Lemma~\ref{lem:rank-condition} gives
\[
        R_r(D) \mid n,
\]
which proves (iv).

Since $R_r(D)$ is a positive integer dividing $n$, and $n\le p_r+1$, any 
supporting branch must satisfy
\[
        R_r(D) \le p_r+1.
\]
This proves the final assertion.
\end{proof}

\begin{corollary}[Branch elimination criterion]\label{cor:branch-elimination}
Let $r \ge 38$ and let $D \mid P_r$ be squarefree.  If
\[
        R_r(D)>p_r+1,
\]
then no prime-complete product of order $r$ has squarefree branch $D$.
\end{corollary}

\begin{proof}
This is the contrapositive of the final assertion of Theorem~\ref{thm:reduction}.
\end{proof}

\section{Interpretation}

Theorem~\ref{thm:reduction} does not assert that the small-index core is empty.  
It asserts that every tail counterexample must lie inside it.  More explicitly, 
a prime-complete product of order $r \ge 38$ can exist only if there is a 
squarefree branch $D\mid P_r$ and an index
\[
        n\le p_r+1
\]
such that
\[
        R_r(D)\mid n,
        \qquad
        B_r(D)\mid Y_n,
        \qquad
        P^+(Y_n)\le p_r.
\]
The rank condition is necessary for the complementary support, and the primitive-divisor 
theorem makes the index small.  Thus the original tail problem is reduced to a restricted 
finite-index obstruction on Pell--Lucas branches.

This reduction is independent of the $abc$ conjecture.  It also avoids the stronger 
problem of bounding the largest pair of consecutive $p_r$-smooth integers.  It uses the 
prime-complete hypothesis directly: every prime up to $p_r$ must occur, and every prime 
not occurring in the squarefree part $D$ must occur in the Pell $Y$-coordinate.

\begin{remark}[Why the theorem is tail-stated]
The proof of the index cap only needs $r \ge 10$, but the theorem is stated for $r \ge 38$ 
because that is the intended tail range after the finite computational base.  
Nothing in the proof depends on the specific value $38$ except the inequality $r \ge 10$.
\end{remark}

\begin{remark}[On possible strengthenings]
For a fixed branch $D$, stronger lower bounds for the greatest prime factor of $Y_n$ may 
eventually improve the cap $p_r+1$ for that branch.  Uniform improvement across all 
$D\mid P_r$ is a separate and more delicate issue because the quadratic field varies 
with $D$.  The reduction theorem above intentionally avoids any such uniformity assumption.
\end{remark}

\begin{thebibliography}{9}

\bibitem{BHV}
Y. Bilu, G. Hanrot, and P. M. Voutier,
\emph{Existence of primitive divisors of Lucas and Lehmer numbers},
J. Reine Angew. Math. \textbf{539} (2001), 75--122.

\bibitem{Lehmer}
D. H. Lehmer,
\emph{On a problem of St\o rmer},
Illinois J. Math. \textbf{8} (1964), 57--79.

\bibitem{LucaNajman}
F. Luca and F. Najman,
\emph{On the largest prime factor of $x^2-1$},
arXiv:1005.1533.

\bibitem{Najman}
F. Najman,
\emph{Smooth values of quadratic polynomials},
arXiv:1005.1535.

\end{thebibliography}

\end{document}
